%======================================================================
\section{Conjectures and open problems}\label{sec:conjectures}
%======================================================================

For each conjecture, Table~\ref{tab:falsify} gives what it would imply and through which result, the evidence, and finite
computations that would refute it; these are sufficient tests, not a complete list of the ways in which the conjecture could fail.

\begin{table}[!ht]
\centering\small
\caption{Conjectures of the paper: consequence, evidence, and sufficient finite refutation tests.}\label{tab:falsify}
\begin{tabular}{@{}>{\raggedright\arraybackslash}p{0.15\textwidth}>{\raggedright\arraybackslash}p{0.22\textwidth}>{\raggedright\arraybackslash}p{0.27\textwidth}>{\raggedright\arraybackslash}p{0.28\textwidth}@{}}
\toprule
Statement & Gives (via) & Evidence & Refuted, for example, by \\
\midrule
Conj.~\ref{conj:U} & $\RH\Leftrightarrow\condE$ (Thm.~\ref{thm:logic}) & Section~\ref{sec:uniqueness}; Cor.~\ref{cor:nearrigid} & a second admissible pair \\
Conj.~\ref{conj:S} & $\kstar=0$ under \condE\ (Thm.~\ref{thm:duality}(a)) & Table~\ref{tab:kappa-ladder} & an admissible pair at $q=1$ with $\mu\ge\lambda\dd t$,
$\lambda>0$; the finite certificates here cannot confirm it \\
Conj.~\ref{conj:family} & \condS, \condU\ (Thm.~\ref{thm:criterion}, Cor.~\ref{cor:magic}(a)) & Props.~\ref{thm:finiteJ} and~\ref{prop:exactcone}; Obs.~\ref{obs:family} & a singular $M_J$;
$p_k^{(J)}\le0$ with $J\ge10$; $a_J\ge\pi/2$; no finite prefix of the family alone refutes (c) \\
\bottomrule
\end{tabular}
\end{table}

\subsection{A prime-power magic-function conjecture}

Let $P_J$, $H_J=P_J(t^2)$ and $F_J=\Xi^2H_J$ be the exact family of Proposition~\ref{prop:E}. The following conjecture is of the same kind as
the conjectures of Cohn and Miller on explicit sequences converging to magic functions \cite{CM16}.

\begin{conjecture}[Prime-power magic function]\label{conj:family}
\begin{enumerate}[label=\textup{(\alph*)}]
\item $M_J$ is non-singular for every $J\ge1$.
\item For every $J\ge10$ all coefficients $p^{(J)}_k$, $0\le k\le2J$, are positive, and
$a^*:=\sup_{J\ge10}\max_{1\le k\le2J}\bigl((2k)!\,p^{(J)}_k\bigr)^{1/(2k)}<\pi/2$.
\item $P_J\to P_\infty$ coefficientwise, and $F_\infty:=\Xi^2P_\infty(t^2)$ satisfies $\FT F_\infty(\xi)\ge0$ for every $\xi\ge\xin2$.
\end{enumerate}
\end{conjecture}

\emph{Implication.} The conjecture implies \condS\ and \condU\ unconditionally, hence $\RH\Leftrightarrow\condE$. Indeed (b) is
\eqref{eq:POSa} with $a=a^*$ and $C=1$, so Proposition~\ref{prop:coeff} gives $H_J\ge1$ on $\RR$ and $|H_J(t)|\le\cosh(a^*|t|)$; with this
bound, coefficientwise convergence gives $H_J\to H_\infty$ locally uniformly with a uniform exponential majorant of type $a^*<\pi/2$, so
$\FT F_J\to\FT F_\infty$ pointwise by dominated convergence, and (c) gives \hyp{Mg}. Remark~\ref{rem:variant}(a) applies. Part (c) is
weaker than in an earlier version of this paper, which asked for $\FT F_\infty>0$ off $\xiPP$; the weakening rests on
Theorem~\ref{thm:zerosupport}.

\emph{Evidence.} Parts (a) and (b) are certified at $J\in\{10,60,61,110,111\}$ (Proposition~\ref{thm:finiteJ}) and computed for
$1\le J\le110$ and $10\le J\le110$ respectively (Numerical observation~\ref{obs:family}), which also shows why (b) starts at $J=10$. At the
same five $J$, the finite-$J$ form of the inequality in (c), $\FT F_J\ge0$ on $[\xin2,\infty)$, is certified (Proposition~\ref{prop:exactcone}).

\emph{Refutation.} Each of the following, at some finite $J$, refutes the conjecture: a singular $M_J$; a non-positive coefficient with
$J\ge10$; $a_J:=\max_{1\le k\le2J}((2k)!\,|p^{(J)}_k|)^{1/2k}\ge\pi/2$. These are sufficient tests, not a complete list. Part (b) asks for
$\sup_{J\ge10}a_J<\pi/2$, which fails if $a_J\to\pi/2$, even when $a_J<\pi/2$ for every $J$. Part (c) concerns the limit: no finite prefix of
the computed family alone refutes it, although a finite computation combined with rigorous bounds for the rest of the family could.

\subsection{Sharpness and uniqueness}

The certificates of Section~\ref{sec:certified} bound $\kstar$ above by positive numbers, so no finite collection of them proves \condS; a
proof needs more, for instance an exact magic function (Corollary~\ref{cor:magic}, Theorem~\ref{thm:criterion}), which a finite
computer-assisted argument for a suitable uniform statement might supply. A finite computation can refute \condS: an explicit admissible pair at
$q=1$ whose zero measure is certified to satisfy $\mu\ge\lambda\dd t$ with $\lambda>0$ would give $\kstar\ge\lambda>0$ by
Lemma~\ref{lem:weak-duality}. Proposition~\ref{thm:notcritical}
certifies pairs of this kind for other archimedean data, at $\xi_2$ and at their natural gaps. Under \condE, \condS\ says $\kstar=0$, that is, every
admissible $\mu$ is purely atomic on one discrete set (Theorem~\ref{thm:crit}), and the optimised degree-one explicit-formula conductor
bound is exactly $1$. Conjecture~\ref{conj:U} would be refuted by a single admissible pair other than $\pz$. The results of
Section~\ref{sec:uniqueness} exclude such pairs when the zero measure is carried by $\Zz\cup\{0\}$ or is purely atomic with integer weights,
whatever the prime measure (Theorems~\ref{thm:zerosupport} and~\ref{thm:integer}), and when it is carried by $\Zz$ and countably many further
points satisfying \eqref{eq:summable} with the prime measure on $\BRS$ (Theorem~\ref{thm:Gfin}). Corollary~\ref{cor:nearrigid} bounds the zero mass off $\Zz$ at low height for every
admissible pair.

\subsection{Open problems}

\begin{enumerate}[label=\textup{(\arabic*)}]
\item\label{op:fractional} \emph{Non-integer weights.} Conjecture~\ref{conj:U} for zero measures with an atom of non-integer weight or a
non-atomic part. For integer weights it is Theorem~\ref{thm:integer}, which settles the Beurling case asked in an earlier version of this
paper. If \condS\ fails, every admissible pair at the conductor $\qmin<1$ has a non-integer weight (Corollary~\ref{cor:integer}).
\item \emph{Extra primes.} Prime-side mass off $\BRS$ together with extra zeros. Without extra zeros, that is, when the zero measure is
carried by $\Zz\cup\{0\}$, every prime-side mass off $\zeta$'s nodes is excluded by Theorem~\ref{thm:zerosupport}; with extra zeros the
problem is open when some zero weight is not an integer (Theorem~\ref{thm:integer}), in particular for atoms at positions $\eta$ with
$e^{4\pi\eta}$ irrational (Remark~\ref{rem:extraprimes}).
\item \emph{Infinitely many extra zeros} on the zero side without a summability condition. The case of countable $E$ with
$\sum_{e\in E}\mu(\{e\})\,|\zeta(\frac12+ie)|\,(1+|e|)^{-1/4+\delta}<\infty$ for some $\delta>0$, asked in an earlier version of this paper, is
settled by Theorem~\ref{thm:Gfin}, with the prime measure on $\BRS$, under the weaker condition \eqref{eq:summable}. Countable $E$ without such a
condition, unless all weights are integers (Theorem~\ref{thm:integer}), and extra zero-side mass that is not atomic, remain open.
\item\label{op:strict-compl} \emph{Strict complementarity.} Does \condU\ imply the existence of an exact magic function
without spurious zeros?
\item \emph{The existence statement} \condE, that is $\kstar\ge0$, without assuming RH: Odlyzko's question
whether explicit-formula bounds proved under GRH hold unconditionally \cite{Odl90}, specialised to $\QQ$.
\end{enumerate}
